\documentclass{article}
\usepackage[T1]{fontenc}
\usepackage{lmodern}
\usepackage{graphicx}
\usepackage{amsmath,amssymb}
\usepackage[a4paper,margin=2cm]{geometry}
\usepackage[absolute,overlay]{textpos}
\setlength{\TPHorizModule}{1mm}
\setlength{\TPVertModule}{1mm}
\usepackage[indonesian]{babel}
\usepackage{hyperref}
\setlength{\emergencystretch}{2em}
% unit: Halaman 32

\begin{document}

% === Halaman 32 (python/native) ===

32

Contoh: Misalkan plainteks adalah


sistem dan teknologi informasi itb

Panjang kunci = 6

 Enkripsi:                             (*buang spasi)


sistem


dantek


nologi


inform


asiitb

 Cipherteks: (baca secara vertical kolom per kolom) 

 SDNIAIAONSSNLFITTOOIEEGRTMKIMB 

   (tanpa spasi)  

 SDNI AIAO NSSN LFIT TOOI EEGR TMKI MB (atau blok 4 huruf)

1. Columnar Transposition Cipher

\end{document}
